\section{Resolvents and conditional path laws}\label{sec:dynamics}

Uniform convergence of energies concerns continuous functions, whereas
the speed measure may be singular with respect to round area.  The
Green estimate bridges this difference: it makes a bounded source on
a set of small quantum mass produce a uniformly small potential.
After removing that potential, a resolvent is locally harmonic, so
the compactness theorem for harmonic functions applies.  We first
give this argument for a deterministic sequence of graphs.  The last
subsection applies it in the coupled random environments and proves
Theorem~\ref{thm:main}.

\subsection{A deterministic convergence criterion}

Let $S^2$ carry its round distance $d_0$ and area $\dd A$.  Let $\mu$
be a nonatomic probability measure with full support.  Suppose the
quadratic form
\begin{equation}\label{dyn:continuum-form}
 \cE(u,u)=\frac12\int_{S^2}|\nabla u|^2\,\dd A,
 \qquad u\in C^\infty(S^2),
\end{equation}
is closable in $L^2(\mu)$, and its closure $(\cE,\mathcal F)$ is a
regular conservative Dirichlet form with an associated continuous-path
diffusion $B$.  We always start $B$ with law $\mu$.  These hypotheses
hold for the Liouville area measures used here; the time-change and
Dirichlet-form constructions are treated in
\cite{GRV,Berestycki}, with general form theory in \cite{FOT}.
Our normalization is the explicit factor $1/2$ in
\eqref{dyn:continuum-form}.

For each $n$, let $\mathcal G_n$ be a finite connected graph, allowing loops and
multiple edges, with vertex set $V_n$ and a strictly positive
probability measure $\mu_n$ on $V_n$.  Write
\[
 (H_nu)(v)=\sum_{vw}\bigl(u(v)-u(w)\bigr),
 \qquad
 \en_n(u)=\sum_{\{v,w\}}(u(v)-u(w))^2.
\]
The first sum counts incidences at $v$, and the second counts
unoriented edges with multiplicity.  Loop summands vanish.
Let $Y^n$ be the chain with generator
\begin{equation}\label{dyn:abstract-generator}
 (\mathcal L_nu)(v)=-\frac{(H_nu)(v)}{2\mu_n(v)}.
\end{equation}
It is reversible with respect to $\mu_n$, and its form is
$\cE_n(u,u)=\en_n(u)/2$.  Choose maps $\pi_n:V_n\to S^2$.
Uniform comparison of functions always means
\[
 \|u_n-u\circ\pi_n\|_{\infty,V_n}\longrightarrow0.
\]
No injectivity of $\pi_n$ is required.

\begin{theorem}[From energies and Green bounds to path laws]
\label{dyn:conditional}
In the deterministic setting just defined, assume the following.
\begin{enumerate}[label=\textup{(\roman*)}]
\item $(\pi_n)_*\mu_n\Rightarrow\mu$, and the maximum projected
edge length tends to zero:
\[
 \max_{\{v,w\}\in E(\mathcal G_n)}d_0(\pi_n(v),\pi_n(w))\longrightarrow0.
\]
\item Every uniformly convergent sequence with continuous limit $u$
satisfies
\[
 \liminf_n\en_n(u_n)\ge\int|\nabla u|^2\,\dd A,
\]
where the right side is $+\infty$ unless $u\in W^{1,2}(S^2)$.
Every smooth $u$ has a uniformly convergent recovery sequence with
energy tending to that integral.
\item For every open $U\subset S^2$, sequences that are graph-harmonic
at all vertices in $\pi_n^{-1}(U)$ and have uniformly bounded
supremum and total energy are locally equicontinuous in $U$ under
$\pi_n$, and locally uniformly precompact there.
Here equicontinuity means that, for every compact $K\subset U$,
\[
 \lim_{\delta\downarrow0}\limsup_n
 \sup_{\substack{\pi_n(v),\pi_n(w)\in K\\
                  d_0(\pi_n(v),\pi_n(w))\le\delta}}
       |u_n(v)-u_n(w)|=0,
\]
and precompactness means uniform comparison with a continuous
function on each compact subset after a common subsequence.
\item For the centered solution operator
\[
 H_nG_ng=\mu_n g,\qquad \mu_n(G_ng)=0,
 \qquad \mu_n(g)=0,
\]
there are constants $C<\infty$ and $\rho>0$, independent of $n$, such
that
\begin{equation}\label{dyn:green-assumption}
 \|G_ng\|_\infty\le C\,\mu_n(|g|)^\rho
 \quad\text{if }\|g\|_\infty\le1\text{ and }\mu_n(g)=0.
\end{equation}
\end{enumerate}
Then the laws of $(\pi_n(Y^n_t))_{0\le t\le T}$, with initial law
$\mu_n$, converge in $D([0,T],S^2)$ to the law of
$(B_t)_{0\le t\le T}$, for every finite $T$.  The limit has continuous
paths.  In addition, for every $\alpha>0$ and $f\in C(S^2)$, the
resolvents converge uniformly under $\pi_n$.
\end{theorem}

We prove the theorem through four lemmas.  Their hypotheses throughout
are those of Theorem~\ref{dyn:conditional}.  Notice that full support
and weak convergence of the measures already imply that $\pi_n(V_n)$
becomes dense: otherwise a fixed open ball of positive $\mu$-mass
would be missed along a subsequence.

\begin{lemma}[Continuous Sobolev functions in the closed form]
\label{dyn:smooth-core}
If $u\in C(S^2)\cap W^{1,2}(S^2)$, then $u\in\mathcal F$ and its
closed-form energy is the right side of
\eqref{dyn:continuum-form}.  It has smooth approximations converging
both uniformly and in the form norm.  Two continuous versions of
the same element of $L^2(\mu)$ are equal everywhere.
\end{lemma}

\begin{proof}
In coordinate disks, mollification of a continuous Sobolev function
converges both uniformly on smaller disks and in Sobolev norm.
A smooth partition of unity gives smooth functions $u_j$ on the
sphere with these two convergences globally.  Since $\mu$ is finite,
uniform convergence also gives $L^2(\mu)$ convergence.  Thus $(u_j)$
is Cauchy for $\cE(\cdot,\cdot)+\|\cdot\|_{L^2(\mu)}^2$ and has
limit $u$ in the stipulated closure.  The gradient convergence
identifies its energy.  Polarization gives the same identity for
the bilinear form on continuous Sobolev functions.  Finally a
nonzero continuous difference is bounded away from zero on some
open set, which has positive $\mu$-mass by full support.
\end{proof}

Let $R^n_\alpha=(\alpha-\mathcal L_n)^{-1}$, and let $R_\alpha$ be the
$L^2(\mu)$ resolvent of $(\cE,\mathcal F)$.

\begin{lemma}[Uniform resolvent convergence]\label{dyn:resolvent}
For every $\alpha>0$ and $f\in C(S^2)$, $R_\alpha f$ has a unique
continuous version, which belongs to round $W^{1,2}$, and
\begin{equation}\label{dyn:resolvent-limit}
 \bigl\|R^n_\alpha(f\circ\pi_n)
              -(R_\alpha f)\circ\pi_n\bigr\|_\infty
       \longrightarrow0.
\end{equation}
The same conclusion holds with right-hand sides $f_n$ converging
uniformly under $\pi_n$ to $f$.
\end{lemma}

\begin{proof}
Put $u_n=R^n_\alpha(f\circ\pi_n)$.  The maximum principle and the
resolvent equation give
\begin{equation}\label{dyn:resolvent-bounds}
 \|u_n\|_\infty\le\frac{\|f\|_\infty}{\alpha},
 \qquad
 H_nu_n=\mu_n g_n,
 \qquad
 g_n=2(f\circ\pi_n-\alpha u_n).
\end{equation}
In particular $g_n$ is uniformly bounded, and multiplying the equation
by $u_n$ shows that $\en_n(u_n)$ is uniformly bounded.

We first prove equicontinuity.  Fix a point $x\in S^2$, and choose
disjoint open neighborhoods $U$ of $x$ and $V$ of a different point.
Take $\mu$-null boundaries, with $\mu(V)>0$.  The neighborhood $U$
may be made arbitrarily small in $\mu$-mass because $\mu$ has no
atoms.  Write $U_n=\pi_n^{-1}(U)$ and $V_n'=\pi_n^{-1}(V)$, and set
\[
 q_n=g_n\1_{U_n}
       -\frac{\mu_n(g_n\1_{U_n})}{\mu_n(V_n')}\1_{V_n'}.
\]
For large $n$, the denominator is bounded below.  Hence $q_n$ has
mean zero, bounded supremum, and
$\mu_n(|q_n|)\le C_1\mu_n(U_n)$.  Scaling
\eqref{dyn:green-assumption} yields
\[
 \|G_nq_n\|_\infty\le C_2\mu_n(U_n)^\rho.
\]
Also
\[
 \en_n(G_nq_n)=\mu_n((G_nq_n)q_n)
 \le\|G_nq_n\|_\infty\mu_n(|q_n|).
\]
Thus $u_n-G_nq_n$ has bounded energy and supremum and is harmonic
on $U_n$.  Hypothesis~(iii) makes this difference equicontinuous on
compact subsets of $U$.  Given an oscillation tolerance, first
choose $U$ so that $G_nq_n$ has smaller supremum, and then use
harmonic equicontinuity inside $U$.  A finite cover of $S^2$ proves
equicontinuity of the original sequence.  Boundedness and density of
$\pi_n(V_n)$ now give uniformly convergent subsequences.

Let $u$ be one such limit.  Hypothesis~(ii) and
Lemma~\ref{dyn:smooth-core} show that $u\in\mathcal F\cap C(S^2)$.
The discrete resolvent minimizes
\[
 J_n(v)=\frac12\en_n(v)+\alpha\mu_n(v^2)
                       -2\mu_n((f\circ\pi_n)v).
\]
The energy lower bound, uniform convergence, and weak convergence of
the measures imply
\[
 \cE(u,u)+\alpha\mu(u^2)-2\mu(fu)
       \le\liminf_n J_n(u_n).
\]
For any smooth $v$, take its uniform recovery sequence $v_n$.
Minimality gives $J_n(u_n)\le J_n(v_n)$, and the latter converges to
$\cE(v,v)+\alpha\mu(v^2)-2\mu(fv)$.  Smooth functions are a core
of the closed form by its definition, so $u$ minimizes this
functional over all of $\mathcal F$.  Strict convexity for $\alpha>0$
identifies $u=R_\alpha f$ in $L^2(\mu)$.  Full support makes its
continuous version unique.  Every subsequence has the same limit,
which proves \eqref{dyn:resolvent-limit}.

Finally, $\|R^n_\alpha a-R^n_\alpha b\|_\infty
\le\alpha^{-1}\|a-b\|_\infty$.  This contraction bound proves the
claim for varying right-hand sides.
\end{proof}

\subsection{Controlling small time intervals}

Resolvent convergence integrates the walk over time.  To recover
the actual clock, we next construct functions that separate spatial
sets and whose discrete generators are bounded.  The needed continuum
approximation is uniform, rather than merely in $L^2(\mu)$.

\begin{lemma}[Uniform resolvent approximation]\label{dyn:feller}
For every $f\in C(S^2)$,
\begin{equation}\label{dyn:approximate-identity}
 \|\alpha R_\alpha f-f\|_\infty\longrightarrow0
 \qquad\text{as }\alpha\longrightarrow\infty.
\end{equation}
\end{lemma}

\begin{proof}
First suppose $f$ is smooth, and put $w_\alpha=\alpha R_\alpha f$.
The spectral theorem for the nonnegative self-adjoint form operator
gives $w_\alpha\to f$ in the form norm.  Indeed the spectral
multiplier $\alpha/(\alpha+\lambda)$ tends to one and is bounded by
one, so dominated convergence applies with spectral weight
$1+\lambda$.  By Lemma~\ref{dyn:smooth-core},
\[
 \|\nabla(w_\alpha-f)\|_{L^2(\dd A)}\longrightarrow0.
\]
The maximum principle, or its discrete version followed by
Lemma~\ref{dyn:resolvent}, gives
\begin{equation}\label{dyn:range}
 \min f\le w_\alpha\le\max f.
\end{equation}

We also need convergence in round $L^2$, which does not follow by
changing the measure in the form norm.  Let $c_\alpha$ be the round
average of $w_\alpha-f$.  Round Poincare gives
\[
 \|w_\alpha-f-c_\alpha\|_{L^2(\dd A)}\longrightarrow0.
\]
The constants $c_\alpha$ are bounded by \eqref{dyn:range}.  If a
subsequence has limit $c$, then $w_\alpha\to f+c$ in round $L^2$
on that subsequence, and therefore along a further subsequence almost
everywhere.  Consequently
$\min f\le f+c\le\max f$ almost everywhere for round area.
Continuity of $f$ makes its essential extrema equal to its extrema:
every neighborhood of a minimizer or maximizer has positive round
area.  The two inequalities force $c\ge0$ and $c\le0$.
Thus $c=0$, and $w_\alpha-f\to0$ in round $W^{1,2}$.

Fix a small coordinate disk about a point $x$, and a fixed annulus
inside that disk surrounding a smaller neighborhood of $x$.
Radial averaging of the round $W^{1,2}$ error supplies, for each
large $\alpha$, a circle in this annulus on which the trace of
$w_\alpha-f$ tends to zero in one-dimensional $H^1$.  The radii stay
in a fixed compact subinterval of $(0,\infty)$, so the
one-dimensional Sobolev inequality gives uniform convergence on
these circles.  Almost every radius has the Sobolev trace of the
continuous representative, which is the trace we choose.

Let $D_\alpha$ be the disk inside the chosen circle.  The resolvent
equation is
\begin{equation}\label{dyn:weak-resolvent}
 \cE(w_\alpha,v)+\alpha\int(w_\alpha-f)v\,\dd\mu=0,
 \qquad v\in\mathcal F.
\end{equation}
Choose a constant $K$ strictly above the boundary values of
$w_\alpha$ and at least $\sup_{D_\alpha} f$.
The positive part $(w_\alpha-K)^+$ on $D_\alpha$, extended by zero,
vanishes in a collar of the circle and is continuous Sobolev.
Lemma~\ref{dyn:smooth-core} makes it an admissible test in
\eqref{dyn:weak-resolvent}.  The two terms are nonnegative; the
second is at least $\alpha\int ((w_\alpha-K)^+)^2\,\dd\mu$.
Thus $w_\alpha\le K$ on the disk, first $\mu$-almost everywhere and
then everywhere by continuity and full support.  The negative-part
test gives the lower bound.  Letting the strict slack tend to zero
shows that the values inside the circle are bounded by its boundary
values and the extrema of $f$ in the disk.

The oscillation of $f$ on sufficiently small disks is arbitrarily
small.  The boundary convergence and this comparison therefore
give uniform convergence on a neighborhood of $x$.  A finite cover
proves \eqref{dyn:approximate-identity} for smooth $f$.
For continuous $f$, approximate it uniformly by a smooth function
and use the contraction $\|\alpha R_\alpha\|_{\infty\to\infty}\le1$.
\end{proof}

\begin{lemma}[An exit estimate uniform over starting vertices]
\label{dyn:tightness}
For $d>0$ and $v\in V_n$, let
\[
 \tau^n_d(v)=\inf\{t\ge0:
   d_0(\pi_n(Y^n_t),\pi_n(v))\ge d\}.
\]
Then
\begin{equation}\label{dyn:exit}
 \lim_{s\downarrow0}\limsup_{n\to\infty}
       \sup_{v\in V_n}\P_v[\tau^n_d(v)\le s]=0.
\end{equation}
The stationary projected path laws are tight in $D([0,T],S^2)$,
and all their limits are supported on continuous paths.
\end{lemma}

\begin{proof}
Cover the sphere by finitely many round disks of radius less than
$d/8$.  For each choose a smooth function $f$ in $[0,1]$, zero on
that disk, and one outside the concentric disk of radius $d/2$.
If $\pi_n(v)$ is in the smaller disk and displacement from
$\pi_n(v)$ reaches $d$, the value of $f$ at the new location is one.

Fix $b>0$.  Lemmas~\ref{dyn:feller} and \ref{dyn:resolvent}, first
with a sufficiently large fixed $\alpha$ and then with sufficiently
large $n$, give for each of this finite family a function
\[
 v_n=\alpha R^n_\alpha(f\circ\pi_n),\qquad
 \|v_n-f\circ\pi_n\|_\infty\le b.
\]
Its generator satisfies
\[
 \mathcal L_nv_n=\alpha(v_n-f\circ\pi_n),
 \qquad \|\mathcal L_nv_n\|_\infty\le C_b,
\]
with a common $C_b$ for the finite family and all sufficiently large
$n$.  Start the chain at $v$ in the corresponding small disk and
stop at $\sigma=\tau^n_d(v)\wedge s$.  If
$p=\P_v[\tau^n_d(v)\le s]$, then
\[
 \E_v[v_n(Y^n_\sigma)]\ge p-b,
 \qquad
 \E_v[v_n(Y^n_\sigma)]
   =v_n(v)+\E_v\int_0^\sigma \mathcal L_nv_n(Y^n_t)\,\dd t
   \le b+C_bs.
\]
Hence $p\le2b+C_bs$, uniformly over all starting vertices.
First let $s\downarrow0$ and then $b\downarrow0$ to obtain
\eqref{dyn:exit}.

Apply the strong Markov property to the chain $Y^n$ after any
stopping time.  Since \eqref{dyn:exit} is uniform over its current
vertex, it gives the usual stopping-time increment condition for
the projected paths.  The projection itself need not be Markov.
Apply Aldous's stopping-time criterion \cite{Aldous} to the three
coordinate functions of the round sphere in $\R^3$; their values
are bounded.  This gives tightness of each coordinate path law.
The largest coordinate jump is bounded by the maximum projected
edge length, which tends to zero.  Under $J_1$ convergence any jump
of a limit is approximated by jumps of the approximating paths,
so every coordinate limit is continuous.  Coordinate tightness
with continuous limits gives joint tightness with continuous
limits in $\R^3$: on continuous limits, $J_1$ convergence agrees
with uniform convergence, allowing a common time comparison for
the three coordinates.  The limits remain in the closed set
$S^2$, proving the assertion.
\end{proof}

\subsection{Identifying the entire conditional law}

\begin{proof}[Completion of the proof of Theorem~\ref{dyn:conditional}]
Take a subsequence of the stationary projected path laws with a weak
limit.  By Lemma~\ref{dyn:tightness}, diagonal extraction on the
horizons $1,2,\ldots$ gives a consistent continuous-path law on
$[0,\infty)$, denoted by $Q$.  Restrictions are consistent because
the limits are continuous at the restriction times.

Let $f_0,\ldots,f_k\in C(S^2)$ and $\alpha_1,\ldots,\alpha_k>0$.
Integrate the stationary finite-dimensional expectation at times
$0,t_1,t_1+t_2,\ldots,t_1+\cdots+t_k$ against
$\prod_{j=1}^k e^{-\alpha_jt_j}\,\dd t_j$.  The Markov property
and the resolvent identity give exactly
\begin{equation}\label{dyn:laplace}
 \mu_n\!\left[
  (f_0\circ\pi_n)R^n_{\alpha_1}
  \left((f_1\circ\pi_n)R^n_{\alpha_2}
    \left(\cdots(f_{k-1}\circ\pi_n)
                R^n_{\alpha_k}(f_k\circ\pi_n)\cdots\right)\right)
 \right].
\end{equation}
Successive applications of Lemma~\ref{dyn:resolvent}, including its
varying-data assertion, and weak convergence of the measures show
that \eqref{dyn:laplace} converges to
\[
 \mu\!\left[f_0R_{\alpha_1}
       \left(f_1R_{\alpha_2}
          \left(\cdots f_{k-1}R_{\alpha_k}f_k\cdots\right)\right)
       \right].
\]
By the semigroup-resolvent relation for the closed form, this is the
same integrated expectation for the stationary diffusion $B$.

It is also the integrated expectation for $Q$.  On bounded time
boxes this follows from weak path convergence, since evaluation at
fixed times is continuous at continuous paths.  Dominated
convergence then integrates over the gaps.  The contribution of
the complement of a large time box is bounded uniformly by the
tails of the product of exponentials, times
$\prod_j\|f_j\|_\infty$.

The product expectations under $Q$ and under $B$ are bounded
continuous functions of $(t_1,\ldots,t_k)$, including zero gaps.
Uniqueness of multivariate Laplace transforms makes these
functions equal: one may apply one-dimensional uniqueness
successively in each coordinate, and then use continuity to
remove almost-everywhere qualifications.  Product tests determine
probability measures on the compact product $(S^2)^{k+1}$.
Thus $Q$ and the law of $B$ have the same finite-dimensional
distributions, and hence the same law on continuous paths.
Every subsequential limit is the same, proving the theorem.
\end{proof}

\subsection{The clock and the metric-space conclusion}

We now return to the random maps.  The exact relation between the
graph energy and the prescribed walk fixes the clock without any
subsequence-dependent factor.

\begin{lemma}[Exact energy of the accelerated walk]\label{dyn:energy-clock}
For the half-edge walk of Theorem~\ref{thm:main},
\[
 -\langle u,nL_nu\rangle_{L^2(\mu_n)}=\frac12\en_n(u).
\]
Equivalently, $nL_n=-\tfrac12\mu_n^{-1}H_n$, where
$\mu_n^{-1}$ denotes division by vertex mass.
\end{lemma}

\begin{proof}
The generator convention gives $H_n=-\deg\,L_n$.  Since
$\mu_n(v)=\deg(v)/(2n)$,
\[
 -\sum_v\mu_n(v)u(v)nL_nu(v)
   =\frac12\sum_vu(v)H_nu(v)
   =\frac12\sum_{\{v,w\}}(u(v)-u(w))^2.
\]
Loops contribute zero to both sides.  Their degree contributions
are already included in the first identity, so no holding-time
correction is needed.
\end{proof}

\begin{proof}[Proof of Theorem~\ref{thm:main}]\label{dyn:main-proof}
Start with any subsequence.  Retain the random variable $A_n$ of
Proposition~\ref{gr:green-bound} together with all the contour,
spatial, and local and whole-sphere array variables used in
Theorem~\ref{sph:energy} and Proposition~\ref{sph:comparison}.
Its uniform first-moment bound gives tightness, so these variables
admit a joint further extraction.  Apply the sphere results to this
enlarged extraction and use a Skorokhod representation.  On the
resulting probability-one event, their conclusions hold and $A_n$
converges to a finite random variable.  In particular
$\sup_n A_n<\infty$ after discarding finitely many initial terms.

Fix a realization in this probability-one event.  Its limiting
quantum sphere is now a fixed measured conformal sphere.  The
projections $\pi_n$, measures, edge mesh, uniform energy bounds and
recovery, and harmonic compactness satisfy hypotheses~(i)--(iii)
of Theorem~\ref{dyn:conditional}.  The realized bounded $A_n$ gives
hypothesis~(iv).  The Liouville form is exactly
\eqref{dyn:continuum-form} with $\mu=\mu_h$.  By
Lemma~\ref{dyn:energy-clock}, the chain in that deterministic theorem
is $Y^n_t=X^n_{nt}$.  Consequently its conditional projected path
law converges to the conditional law of $B^h$, for every fixed
finite horizon.  All probabilities in this invocation are taken
on the individual fixed graphs.  The random object being compared
is the conditional probability law itself.

It remains to express this comparison in the metric $\delta_T$.
The round topology and the $D_h$ topology agree on the compact
limiting sphere, so the identity map between them is uniformly
continuous in both directions.  The path-law convergence just
proved therefore also holds in $D([0,T],(S^2,D_h))$.

The correspondences in Proposition~\ref{sph:comparison} have
distortions tending to zero.  The standard correspondence gluing
constructs a common metric space containing isometric copies of
all $(V_n,a_nd_n)$ and $(S^2,D_h)$, with
\[
 \sup_{v\in V_n}\dist_Z\bigl(v,\pi_n(v)\bigr)\longrightarrow0.
\]
For completeness, for each $n$ attach its copy to the sphere by
cross-distances obtained by taking the infimum over corresponding
pairs, adding numbers $\varepsilon_n\downarrow0$ larger than half
the respective distortions;
use the shortest-path metric on the union, identifying only the
common sphere, and then take the completion.  This preserves the
metrics of every copy.  The images become close in Hausdorff
distance, and their measures become close in Prokhorov distance.

Apply the embedding to each coordinate of a chain path.  Its
distance in the common Skorokhod space from its $\pi_n$-projected
path is bounded by the displayed supremum, using the identity time
change.  The corresponding probability laws therefore have
Prokhorov distance tending to zero.  Combining this with the
projected law convergence shows that all three terms defining
$\delta_T$ tend to zero in the coupled extraction.

We have obtained almost-sure convergence in a further coupling from
every original subsequence, to the prescribed quantum sphere and
its conditional diffusion law.  The subsequence principle gives
convergence in distribution through all positive integers $n$.
Forgetting the root, the spanning tree, and the contour marks
leaves precisely the states in Theorem~\ref{thm:main}.  The same
choice $c=1$ works for every time horizon.
\end{proof}
